\section{新形式理论（Atkin--Lehner）}
\label{sec:newforms}
% ============================================================================

现在我们终于碰到 Hecke 理论里最容易被忽略、但在算术上最关键的一层细化。前几节已经知道：在给定级别的空间里，可以找 Hecke 算子的公共特征向量。可是一旦级别从 $M$ 升到 $N$，旧问题马上出现：若 $f \in \Sk_k(M)$，而 $M \mid N$，那么 $f(d\tau)$ 往往自动落在 $\Sk_k(N)$ 里。换句话说，高级别空间里混进了许多从低级别\textbf{抬升}上来的形式。若不把这些"旧信息"剥掉，我们就无法回答最核心的问题：\emph{到底哪些特征形式真正属于级别 $N$，而不是在更低级别就已经出现过？} 新形式理论的核心想法正是把 $\Sk_k(N)$ 分成两块：一块是从低级别搬上来的旧形式，另一块是与它们正交、真正第一次在级别 $N$ 出现的新形式。这样做的收获极大：新形式不仅给出最干净的特征形式基，而且会与椭圆曲线、$L$-函数以及费马大定理的证明链直接对接。

先看一个最具体的例子。设
\[
  \Delta(\tau)=q-24q^2+252q^3-1472q^4+\cdots \in \Sk_{12}(\SL_2(\Z)).
\]
把变量换成 $2\tau$，得到
\[
  \Delta(2\tau)=q^2-24q^4+252q^6-1472q^8+\cdots.
\]
它不再是级别 $1$ 的形式，却是级别 $2$ 的尖点形式。也就是说，$\Sk_{12}(\GammaO(2))$ 里至少有一部分内容，其实早在级别 $1$ 就已经存在了，只是被重新嵌进了更大的空间。这个例子说明：在讨论级别 $N$ 的 Hecke 特征形式时，单看“是不是特征向量”还不够，我们还必须问一句：\textbf{它是不是第一次在这个级别出现？}

\begin{definition}[退化映射与旧形式空间]
\label{def:oldforms}
设 $M \mid N$，$d \mid N/M$。定义\textbf{退化映射}（degeneracy map）
\[
  V_d\colon \Sk_k(M) \longrightarrow \Sk_k(N),\qquad (V_df)(\tau)=f(d\tau).
\]
于是级别 $N$ 的\textbf{旧形式空间}（oldspace）\index{旧形式 (oldform)}定义为
\[
  \Sk_k^{\mathrm{old}}(N)
  = \sum_{\substack{M\mid N \\ M<N}}\ \sum_{d\mid N/M} V_d\Sk_k(M)
  = \sum_{\substack{M\mid N \\ M<N}}\ \sum_{d\mid N/M}
    \{f(d\tau):f\in \Sk_k(M)\}.
\]
\end{definition}

\begin{example}[级别 $2$ 里的旧形式]
上面的 $\Delta(2\tau)$ 正是 \cref{def:oldforms} 里的情形：这里 $M=1$，$N=2$，$d=2$，所以
\[
  \Delta(2\tau)=V_2\Delta\in \Sk_{12}^{\mathrm{old}}(2).
\]
它的 Fourier 展开从 $q^2$ 开始，清楚地反映出“把同一个级别 $1$ 的对象沿着 $\tau\mapsto 2\tau$ 搬到级别 $2$”这一过程。
\end{example}

要定义新形式，就需要把 Petersson 内积从 \cref{def:petersson} 的级别 $1$ 情形推广到 $\GammaO(N)$：
若 $\mathcal{F}_N$ 是 $\GammaO(N)$ 的一个基本域，令
\[
  \langle f,g\rangle_N
  =\int_{\mathcal{F}_N} f(\tau)\overline{g(\tau)}\,(\Im\tau)^k\,\frac{dx\,dy}{y^2}.
\]
由于尖点形式在各尖点处都指数衰减，积分仍然收敛。

\begin{definition}[新形式空间]
\label{def:newforms-space}
级别 $N$ 的\textbf{新形式空间}（newspace）\index{新形式 (newform)}定义为旧形式空间在 Petersson 内积下的正交补：
\[
  \Sk_k^{\mathrm{new}}(N)
  = \bigl(\Sk_k^{\mathrm{old}}(N)\bigr)^{\perp}
  \subset \Sk_k(N).
\]
位于 $\Sk_k^{\mathrm{new}}(N)$ 中、首项归一化且是所有相关算子的公共特征向量的形式，称为\textbf{新形式}。
\end{definition}

% TikZ figure: oldforms vs newforms
% Shows the inclusion/projection between levels N and M|N
\begin{figure}[ht]
\centering
\begin{tikzpicture}[
  box/.style={draw, rounded corners=4pt, minimum width=2.8cm, minimum height=1cm, font=\small, align=center},
  arrow/.style={->, thick, >=stealth}
]

% Level M box
\node[box, fill=green!15] (SM) at (0, 3) {$\Sk_k(M)$\\(低级别)};

% Map down to level N
\node[box, fill=blue!10] (SN) at (4, 3) {$\Sk_k(N)$\\(高级别 $M \mid N$)};

% Oldforms inside SN
\node[box, fill=orange!15, dashed, draw=orange] (old) at (4, 1) {旧形式 $\Sk_k^{\mathrm{old}}(N)$};

% Newforms inside SN
\node[box, fill=red!10, dashed, draw=red] (new) at (7.5, 3) {新形式 $\Sk_k^{\mathrm{new}}(N)$};

% Arrows
\draw[arrow] (SM) -- node[above, font=\small] {$f(\tau) \mapsto f(d\tau)$, $d = N/M$} (SN);
\draw[arrow, dashed] (SN.south) -- (old.north) node[midway, right, font=\small] {投影};

% Decomposition
\draw[arrow, double] (SN.east) -- node[above, font=\small] {正交分解} (new.west);

% Horizontal brace for decomposition
\node[font=\small] at (5.5, 4.0) {$\Sk_k(N) = \Sk_k^{\mathrm{old}}(N) \oplus \Sk_k^{\mathrm{new}}(N)$};

% Annotations
\node[font=\scriptsize, text=orange!80!black, align=center] at (4, 0.1)
  {由低级别 $M \mid N, M < N$ 的形式\\经 $f \mapsto f(d\tau)$ 生成};
\node[font=\scriptsize, text=red!80!black, align=center] at (7.5, 2)
  {与所有旧形式正交\\满足强重数一定理};

\end{tikzpicture}
\caption{$\Sk_k(N)$ 的旧形式—新形式正交分解。旧形式由低级别升上来，新形式是"真正属于级别 $N$"的部分。}
\label{fig:oldforms-newforms}
\end{figure}


\begin{theorem}[旧空间与新空间的正交分解]
\label{thm:newform-decomposition}
对任意整数 $N\ge 1$，有正交直和分解
\[
  \Sk_k(N)=\Sk_k^{\mathrm{old}}(N)\oplus \Sk_k^{\mathrm{new}}(N).
\]
\end{theorem}

\begin{proof}
空间 $\Sk_k(N)$ 是有限维复向量空间；给定 Petersson 内积后，它成为有限维 Hermite 空间。任取其子空间
$U=\Sk_k^{\mathrm{old}}(N)$，线性代数告诉我们：正交补
\[
  U^{\perp}=\{h\in \Sk_k(N):\langle h,u\rangle_N=0\text{ 对所有 }u\in U\}
\]
也是子空间，并且满足 $\Sk_k(N)=U\oplus U^{\perp}$。

这里的直和是正交直和：若 $x\in U\cap U^{\perp}$，则 $\langle x,x\rangle_N=0$，由内积正定性得到 $x=0$。而每个 $f\in \Sk_k(N)$ 都可唯一写成
$f=u+h$，其中 $u\in U$，$h\in U^{\perp}$。按照 \cref{def:newforms-space}，$U^{\perp}=\Sk_k^{\mathrm{new}}(N)$，命题得证。
\end{proof}

\begin{example}[级别 $1$ 时没有旧形式]
\label{ex:level1-all-new}
当 $N=1$ 时，不存在满足 $M\mid 1$ 且 $M<1$ 的正整数 $M$，所以
\[
  \Sk_k^{\mathrm{old}}(1)=0,
  \qquad
  \Sk_k^{\mathrm{new}}(1)=\Sk_k(1).
\]
因此级别 $1$ 的所有尖点形式自动都是新形式。特别地，
$\dim \Sk_{12}(\SL_2(\Z))=1$，故 $\Delta$ 不只是 Hecke 特征形式，而且是重量 $12$、级别 $1$ 的唯一归一化新形式。
\end{example}

接下来需要一个能看见“坏素数处局部对称性”的算子族。Hecke 算子 $T_p$（$p\nmid N$）控制的是好素数处的信息；而整除 $N$ 的素数还携带额外的局部对称，这正由 Atkin--Lehner 算子来记录。

\begin{definition}[恰除因子与 Atkin--Lehner 算子]
\label{def:atkin-lehner}
若 $Q\mid N$ 且 $\gcd(Q,N/Q)=1$，则称 $Q$ 是 $N$ 的\textbf{恰除因子}（exact divisor），记作
$Q\mid\mid N$\index{恰除因子 (exact divisor)}。在这种情形下，可取整数 $a,b,c,d$ 使得
\[
  Qad-\frac{N}{Q}bc=1,
\]
于是矩阵
\[
  w_Q=
  \begin{pmatrix}
    Qa & b\\
    Nc & Qd
  \end{pmatrix}
\]
的行列式是 $Q$。对 $f\in \Sk_k(N)$，定义 \textbf{Atkin--Lehner 算子}\index{Atkin--Lehner 算子}
\[
  W_Qf = Q^{-k/2}\,(f\mid_k w_Q).
\]
不同的 $w_Q$ 给出同一个作用在 $\Sk_k(N)$ 上的算子；特别地，$W_1=\mathrm{id}$，而 $W_N$ 常称为\textbf{Fricke 对合}。
\end{definition}

\begin{example}[恰除因子的样子]
若 $N=12$，则恰除因子恰好是
\[
  1,\ 3,\ 4,\ 12.
\]
例如 $Q=3$ 是恰除因子，因为 $\gcd(3,4)=1$；但 $Q=2$ 不是，因为
$\gcd(2,6)=2\neq 1$。因此级别 $12$ 上可以讨论 $W_3,W_4,W_{12}$，却不能讨论 $W_2$。
\end{example}

\begin{theorem}[新空间上的同时对角化]
\label{thm:newforms-eigen}
设 $N\ge 1$。则：
\begin{enumerate}
  \item 对每个满足 $\gcd(n,N)=1$ 的整数 $n$，Hecke 算子 $T_n$ 保持
        $\Sk_k^{\mathrm{old}}(N)$ 与 $\Sk_k^{\mathrm{new}}(N)$；
  \item 对每个恰除因子 $Q\mid\mid N$，Atkin--Lehner 算子 $W_Q$ 也保持
        $\Sk_k^{\mathrm{old}}(N)$ 与 $\Sk_k^{\mathrm{new}}(N)$；
  \item 因而 $\Sk_k^{\mathrm{new}}(N)$ 有一组正交基，由同时满足
        \[
          T_n f=a_n(f)f\quad (\gcd(n,N)=1),
          \qquad
          W_Qf=\varepsilon_Q(f)f\quad (Q\mid\mid N)
        \]
        的向量组成，其中每个 $\varepsilon_Q(f)=\pm 1$。
\end{enumerate}
换言之，新形式可以取为所有好 Hecke 算子与所有 Atkin--Lehner 算子的公共特征向量。
\end{theorem}

\begin{proof}
先证 (1)。任取旧形式生成元 $V_df=f(d\tau)$，其中 $f\in \Sk_k(M)$，$M\mid N$，$M<N$，$d\mid N/M$。对任意 $n$ 满足 $\gcd(n,N)=1$，特别有 $\gcd(n,d)=1$。把 \cref{thm:hecke-on-q-exp} 的公式分别作用在 $f$ 与 $V_df$ 的 Fourier 展开上，可直接比较系数得到
\[
  T_n(V_df)=V_d(T_nf).
\]
右边仍在 $V_d\Sk_k(M)$ 中，因此 $T_n$ 保持旧空间。再由 \cref{thm:hecke-selfadjoint} 的同样思路推广到 $\GammaO(N)$ 可知，$T_n$ 关于 Petersson 内积自伴；若 $h\in \Sk_k^{\mathrm{new}}(N)$，$u\in \Sk_k^{\mathrm{old}}(N)$，则
\[
  \langle T_nh,u\rangle_N=\langle h,T_nu\rangle_N=0,
\]
因为 $T_nu\in \Sk_k^{\mathrm{old}}(N)$。所以 $T_nh\in \Sk_k^{\mathrm{new}}(N)$，(1) 成立。

对 (2)，Atkin--Lehner 理论中的标准矩阵计算表明：$W_Q$ 把退化映射的像送到退化映射的像，因此保持旧空间；另一方面 $W_Q$ 是保 Petersson 内积的对合，所以若 $h\perp \Sk_k^{\mathrm{old}}(N)$，则 $W_Qh$ 仍与旧空间正交，于是也落在新空间里。

现在证 (3)。在有限维 Hermite 空间 $\Sk_k^{\mathrm{new}}(N)$ 上，所有 $T_n$（$\gcd(n,N)=1$）彼此交换，所有 $W_Q$ 也彼此交换，并且 $W_Q^2=1$，故每个 $W_Q$ 都可对角化，特征值只能是 $\pm1$。这些算子又都保持新空间，所以可以在新空间上同时对角化，得到一组公共特征向量基。把这些向量再做首项归一化，便得到新形式。
\end{proof}

\begin{example}[级别 $11$ 的唯一新形式]
\label{ex:level11-newform}
空间 $\Sk_2(\GammaO(11))$ 的维数是 $1$，因此它自动等于自己的新空间。故其中任一非零元素都是所有好 Hecke 算子的特征向量；取首项归一化后得到唯一新形式
\[
  f=q-2q^2-q^3+2q^4+q^5+\cdots.
\]
由于维数只有 $1$，$W_{11}$ 也只能按数乘作用，所以 $f$ 还是 $W_{11}$ 的特征向量。这正是新形式理论在最小非平凡级别中的典型样子：一个真正属于级别 $11$ 的对象，把所有算术信息都压缩进一串系数 $a_p(f)$ 里。
\end{example}

这个级别 $11$ 的例子还有更深的几何意义。它对应于导子为 $11$ 的椭圆曲线
\[
  E\colon y^2+y=x^3-x^2-10x-20.
\]
模性定理告诉我们，$E$ 的 Hasse--Weil $L$-函数满足
\[
  L(E,s)=L(f,s).
\]
于是新形式不只是“模形式空间里的漂亮基向量”，而是真正把解析对象与代数几何对象连接起来的桥梁。

\begin{theorem}[强重数一定理]
\label{thm:strong-multiplicity-one}
\index{强重数一定理 (strong multiplicity one)}
设 $f,g\in \Sk_k^{\mathrm{new}}(N)$ 是同一重量、同一级别的归一化新形式。若
\[
  a_p(f)=a_p(g)
\]
对除了有限多个以外的所有素数 $p$ 都成立（等价地，只需对几乎所有 $p\nmid N$ 成立），则
\[
  f=g.
\]
\end{theorem}

\begin{proof}[证明思路]
由 \cref{thm:L-function-euler-product}，$f$ 与 $g$ 的 $L$-函数在几乎所有素数处有相同的局部 Euler 因子，因此它们的商只差一个有限 Euler 乘积。另一方面，归一化新形式满足完整的解析延拓与函数方程；把这些解析性质与 $\GL_2$ 上自守表示的强重数一定理结合起来，就得到对应的自守表示必须同构，从而 $f=g$。

若用现代语言，这一定理是 Jacquet--Langlands 强重数一定理在经典模形式上的具体化；在本书的层次上，可以把它理解为：\textbf{一个新形式由几乎所有 $a_p$ 唯一决定。}
\end{proof}

\begin{example}[为什么这一定理如此强]
若某个归一化新形式 $g$ 满足 $a_p(g)=a_p(\Delta)$ 对除了有限多个以外的所有素数都成立，那么 \cref{thm:strong-multiplicity-one} 立刻推出 $g=\Delta$。也就是说，我们不必逐项比较全部 Fourier 系数；“几乎所有素数处相同”已经足够恢复整个形式。
\end{example}

\subsection*{系数域、Galois 表示与同余数}

CSS 第 II、IV 部分接下来的一个核心思想是：一个新形式并不只是若干复系数的幂级数，它还天然决定了数域、局部域、Galois 表示以及控制同余的 Hecke 局部环。

\begin{definition}[Hecke 系数域]
\label{def:hecke-field}
设 $f=\sum_{n\ge 1}a_n(f)q^n$ 是归一化新形式。定义它的\textbf{Hecke 系数域}
\[
  K_f=\Q\bigl(a_1(f),a_2(f),a_3(f),\ldots\bigr)=\Q\bigl(\{a_n(f)\}_{n\ge 1}\bigr).
\]
这是一个数域；它的整数环记作 $\mathcal{O}_{K_f}$。
\end{definition}

\begin{theorem}[Deligne 定理的预告]
\label{thm:deligne-preview}
设 $f$ 是权 $k$、级别 $N$、角色 $\varepsilon$ 的归一化新形式，$\lambda\mid \ell$ 是 $K_f$ 中一个位于素数 $\ell$ 之上的素理想。则存在连续 Galois 表示
\[
  \rho_{f,\lambda}\colon \Gal(\overline{\Q}/\Q)
  \longrightarrow \GL_2\bigl((K_f)_\lambda\bigr),
\]
满足：对每个 $p\nmid N\ell$，它在 $p$ 处不分歧，而且
\[
  \Tr\bigl(\rho_{f,\lambda}(\Frob_p)\bigr)=a_p(f),
  \qquad
  \det\bigl(\rho_{f,\lambda}(\Frob_p)\bigr)=\varepsilon(p)p^{k-1}.
\]
在本书里，这一结果将在第 9 章作为 Deligne 定理出现。
\end{theorem}

于是，一个新形式给出的不再只是一列 Hecke 特征值，而是一整个相容的 $\lambda$-进表示族。第 9 章的观点正是：模形式的 Fourier 系数、Hecke 代数元素与 Frobenius 迹，实际上是同一份数据的三种语言。

\begin{definition}[同余数]
\label{def:congruence-number}
设
\[
  \phi_f\colon \mathbb{T}_k(N)\longrightarrow \mathcal{O}_{K_f},
  \qquad T_n\longmapsto a_n(f)
\]
是由 $f$ 给出的 Hecke 特征值同态，记 $I_f=\ker(\phi_f)$。把 $\mathrm{Ann}_{\mathbb{T}_k(N)}(I_f)$ 通过 $\phi_f$ 送到 $\mathcal{O}_{K_f}$，得到的理想
\[
  \eta_f=\phi_f\bigl(\mathrm{Ann}_{\mathbb{T}_k(N)}(I_f)\bigr)
\]称为 $f$ 的\textbf{同余理想}；它的绝对范数常称为\textbf{同余数}（congruence number）。
\end{definition}

直观地说，$\eta_f$ 衡量的是：$f$ 的 Hecke 特征值系统与别的模形式“靠得有多近”。若某个素理想 $\lambda\mid \eta_f$，就意味着存在别的特征形式 $g\neq f$ 与 $f$ 在模 $\lambda$ 意义下发生同余。

\begin{example}[Ramanujan 的模 $691$ 同余]
\label{ex:ramanujan-691}
权 $12$、级别 $1$ 的情形给出最著名的例子。由于
\[
  B_{12}=-\frac{691}{2730},
\]
Eisenstein 级数 $E_{12}$ 的分母里出现了素数 $691$。另一方面，$\Sk_{12}(\SL_2(\Z))$ 由 $\Delta$ 张成，于是比较权 $12$ 模形式的 $q$-展开可得经典同余
\[
  \tau(n)\equiv \sigma_{11}(n) \pmod{691}
  \qquad (n\ge 1).
\]
特别地，对每个素数 $p$，
\[
  \tau(p)\equiv 1+p^{11} \pmod{691}.
\]
这说明 cusp 特征值 $\tau(p)$ 与 Eisenstein 特征值 $1+p^{11}$ 在模 $691$ 意义下不可区分。
\end{example}

在 Hecke 代数语言里，这个同余对应于一个 Eisenstein 极大理想
\[
  \mathfrak{m}_{691}=\bigl(691,\ T_p-(1+p^{11})\text{ 对所有 }p\bigr)
  \subset \mathbb{T}_{12}(1).
\]
把 $\mathbb{T}_{12}(1)$ 在 $\mathfrak{m}_{691}$ 处局部化后，我们同时看见 cusp 部分与 Eisenstein 部分被同一个模 $691$ 残差系统“黏”在一起。正是这种局部结构，使 Hecke 代数能够进入 Diamond、Mazur、Wiles 所研究的同余与形变理论。

\textbf{与费马大定理的联系。}
Wiles 证明了每条半稳定有理椭圆曲线都是模的：也就是说，它的 $L$-函数等于某个新形式的 $L(s,f)$。Frey 曲线若真的来自费马方程的非平凡解，那么它是一条半稳定椭圆曲线，因此也应该对应某个适当级别的新形式；但 Ribet 的降级定理表明，这会迫使它对应到一个根本不存在的新形式。矛盾就在这里爆发。换句话说，\textbf{新形式理论把“哪条椭圆曲线来自哪个级别”这个问题组织成了一个足够刚性的框架}，从而 Wiles 与 Ribet 的论证才能真正闭合。


% ============================================================================
